% !TeX root = main.tex
\chapter*{ЗАГАЛЬНІ ВИСНОВКИ}
\addcontentsline{toc}{chapter}{ЗАГАЛЬНІ ВИСНОВКИ}

У дисертації одержано результати відповідно до восьми поставлених завдань щодо формування інформаційної технології виявлення ознак аномального вмісту в мультимедійних даних вебсередовища. Для кожного результату розмежовано формалізований внесок, наявну експериментальну опору та межі допустимого висновку.

\begin{enumerate}[label=\arabic*.]
\item \textit{Формально одержано.} Проаналізовано використання мультимедійних об’єктів як засобів прихованого впливу, способи прихованого розміщення даних і відомі методи виявлення відповідних ознак. Установлено, що ознаки прихованого впливу можуть міститися в декодованому вмісті, частотному або хвильковому представленні, структурі контейнера, службових полях, додаткових байтових областях і подіях оброблення. Статистичні та нейромережеві методи спостерігають слабкі зміни зображення, структурні засоби виявляють невідповідності формату, а подієво-поведінкові засоби реєструють реакцію компонентів вебсистеми. Жодна група окремо не охоплює повний маршрут від носія до можливого впливу. Це обґрунтувало потребу зберігати походження різнорідних свідчень і формувати рішення за їх узгодженістю, а не за механічним максимумом або обов’язковим одночасним спрацюванням. \textit{Внутрішня агрегована опора.} Категоріальні та компонентні підсумки контрольного прогону лише опосередковано узгоджуються з різною реакцією статичних груп ознак; вони не є незалежною перевіркою результатів аналізу джерел. \textit{Не перевірено.} Повноту маршруту від носія до фактичного впливу та вичерпність сукупності методів експериментально не встановлено.

\item \textit{Формалізовано.} Процес прихованої доставки впливу подано через штатні операції вебсередовища. Мультимедійний носій проходить приймання, декодування, перетворення, зберігання та відображення, а можливий прояв пов’язано з конкретною операцією й компонентом оброблення. Відокремлено джерело підготовленого носія, параметри зміни, канал надходження, незмінний вхід і спостережувані наслідки. Така формалізація охоплює випадки, коли носій залишається зовні правдоподібним і формально коректним, але набуває нетипової інтерпретації в окремому засобі оброблення. Визначено, що авторизація джерела, заявлений тип медіаданих або успішне відтворення не скасовують початкового недовіреного статусу носія. \textit{Внутрішня агрегована опора.} В авторських агрегатах наведено результати для сценаріїв послідовної заміни молодших бітів, хвилькового перетворення й хвостових областей; вони характеризують лише контрольовану зміну носія. \textit{Не перевірено.} Виконання аномального коду, його активацію конкретним компонентом вебсистеми та причинно пов’язаний системний наслідок не відтворено.

\item \textit{Формально розроблено.} Модель мультимодального представлення прихованої загрози подає мультимедійний об’єкт як багаторівневий інформаційний контейнер із незмінним байтовим поданням, контрольовано декодованим вмістом, структурною картою, службовим контекстом і окремим причинно пов’язаним журналом подій. Структурне представлення ознак $Z_{s,i}$, спектрально-зорове представлення $Z_{v,i}$ і доступне подієво-поведінкове представлення $Z_{b,i}$ інтегруються у представлення $Z_i$ з явною маскою доступності. Воно зберігає часткові оцінки, локалізації, контекст і посилання на журнал. Модель задає типізовані формати обміну для методу й формалізованої процедури розділу~3 та відокремлює аналітичний вихід від технологічної дії. \textit{Експериментальна опора.} Для внутрішньої серії та чотирьох зовнішніх наборів збережено пооб’єктні ознаки, неперервні оцінки й рішення; зовнішні набори охопили $38$ природних записів із $12$ пристроїв і $24$ груп походження. Ретроспективний прогін має лише агреговані підсумки. \textit{Не перевірено.} Повне інтегроване представлення $Z$ із причинно пов’язаними подіями, переносимість за межі п’яти досліджених форматів і повну подієво-поведінкову конфігурацію не встановлено.

\item \textit{Формально вдосконалено.} Метод спектрально-зорового та структурного аналізу мультимедійних даних поєднує контрольоване декодування, високочастотні залишки, характеристики дискретного косинусного й дискретного хвилькового перетворень та ентропійні характеристики з перевіркою заголовків, меж контейнера, службових полів, завершальних маркерів і додаткових областей. Дві статичні гілки формують окремі оцінки й локалізації, після чого інтегруються зі збереженням позначок доступності. \textit{Експериментальна опора.} У внутрішній серії сформовано $200$ первинних груп і $400$ PNG- та MP4-об’єктів із поділом $60:20:20$. На перевірній частині спільна числова проєкція $Z_c$ дала $40$ правильно визначених контрольних, жодного хибнопозитивного, два хибнонегативні та $38$ правильно визначених змінених об’єктів; частка правильних рішень становила $0{,}9750$, повнота --- $0{,}9500$, $F_1=0{,}9744$, а площа під характеристичною кривою --- $0{,}9869$. Перша зовнішня серія на $56$ похідних MP4 дала збалансовану частку правильних рішень $0{,}8333$ для структурної, $0{,}7857$ для спільної та $0{,}5000$ для спектрально-зорової конфігурації. Багатоформатна серія на $10\,000$ об’єктах перевірила PNG, JPEG, WebP, MP4 і WebM: для структурної, спектрально-зорової та спільної конфігурацій одержано відповідно $0{,}7636$, $0{,}5001$ і $0{,}7916$. Структурна конфігурація мала $0{,}5000$ на WebM, тоді як значення спільної проєкції для п’яти форматів перебували в межах від $0{,}7387$ до $0{,}8173$. Недеструктивне повторне виділення ознак із тих самих $56$ і $2000$ фізичних MP4 зберегло збалансовану правильність структурної гілки $0{,}8333$, але змінило значення спільної проєкції від $0{,}7857$ до $0{,}6905$ і від $0{,}7987$ до $0{,}6520$. Розвідувальний аналіз $1500$ допустимо змінених контрольних об’єктів показав частку хибнопозитивних структурних рішень $1{,}0000$ для приватних розширень або відступів у кожному форматі; усі об’єкти пройшли розбір і декодування, а $1000$ контейнерних пар були піксельно тотожними. У розвідувальному парному аналізі всіх $2500$ форматних пар спектральна умова після декодування завжди мала ненульову відмінність, а середній за відеоджерелами PSNR становив від $38{,}6417$ до $41{,}4147$~дБ; усі $5000$ контейнерних контрольних пар були піксельно тотожними. Отже, встановлено форматну межу й недостатню специфічність структурного складника щодо допустимої контейнерної варіативності, підтверджено його відтворюваність на незмінених фізичних об’єктах і відкинуто повне зникнення спектрального сліду як єдине пояснення низької якості $Z_v$; зовнішньої роздільної здатності та числової стабільності ізольованих спектрально-зорових ознак не підтверджено. \textit{Не перевірено.} Не оцінено інші формати, всі варіанти кодеків і профілів, специфічність після нового навчання з репрезентативними допустимими структурами, збереження первинної просторово-часової локалізації, повну модель зорового перетворювача та виявлення виконуваного шкідливого коду.

\item \textit{Формалізовано.} Процедура поведінкового узгодження ознак різного походження задає зіставлення статичного опису з подіями оброблення того самого об’єкта за ідентифікатором, часовими межами, фазою та причинним зв’язком. Визначено кодування подій, фазову маску, кероване об’єднання, оцінку повноти й посилання на незмінний журнал для контрольної перевірки. Відсутність журналу позначається як недоступність модальності та не трактується як підтвердження норми. \textit{Експериментальна опора.} У ретроспективному прогоні вилучення скороченої правилової оцінки не змінило агрегованих показників, а доступна технічна ознака не спрацювала у $4832$ унікальних відеоматеріалах. Окремо сформовано скорочену модель узгодження, яка враховує три статичні оцінки першого методу, дві їхні розбіжності та ознаку формату. Ця модель узгоджує вже сформовані статичні свідчення і не є реалізацією повної процедури поведінкового узгодження. \textit{Не перевірено.} Причинно пов’язаних журналів $X_b$ не сформовано, тому кодувальник подій, перехресне зіставлення та перевагу повної конфігурації процедури не оцінено.

\item \textit{Сформовано й локально перевірено.} Інформаційна технологія виявлення ознак аномального вмісту в мультимедійних даних вебсередовища передбачає рівно п’ять послідовних етапів: попереднє оброблення, формування ознак, інтеграцію ознак, прийняття рішення та реєстрацію результатів. Наукова відмінність полягає не в числі етапів, а в інваріантах: початковий об’єкт зберігається незмінним; реконструйований об’єкт отримує новий ідентифікатор і проходить окремий цикл; аналітичне рішення відокремлено від дії політики; схема контрольного запису містить контрольну суму входу, версії складників і політики, ідентифікатор запуску, вихідний стан, маску доступності, локалізації, посилання на журнал і часові межі. \textit{Експериментальна опора.} За наперед зафіксованим протоколом виконано $45$ сценарних спостережень для десяти об’єктів п’яти форматів. Створено $35$ цілісних записів $S_5$; п’ять помилок декодування і п’ять ін’єкованих відмов реєстрації завершилися без $S_5$ та без зовнішньої дії. Перевірено порядок $S_1$--$S_5$, незмінність входу, $m_b=0$ за відсутності $J_b$, атомарну реєстрацію без перезапису, повтор, відновлення й один чотирипроцесний запуск. Медіана послідовного циклу становила $20{,}321$~мс, 95-й процентиль --- $108{,}949$~мс; локальна пропускна здатність чотирипроцесного запуску дорівнювала $22{,}827$ об’єкта за секунду за найбільшої сукупної пам’яті $545{,}52$~МіБ. \textit{Не перевірено.} Реконструкцію, причинно пов’язаний журнал, виконання зовнішніх приписів, розподілену чергу, мережеве передавання, багатовузлове сховище й тривале промислове навантаження не оцінено.

\item \textit{Формально описано.} Визначено прототипну конфігурацію з контрольного опису прогону від 11 липня 2026~р. та формат експериментального запису, потрібного для повторної перевірки. Запис кожного об’єкта має поєднувати походження даних, фактично використані представлення, неперервні оцінки, рішення й часові значення; пакет запуску додатково фіксує редакцію програмного коду, перелік залежностей, конфігурацію, початкові значення генераторів, середовище та час фіксації порога. \textit{Експериментальна опора.} Для чотирьох зовнішніх наборів і окремого ресурсного випробування збережено програмний код, реєстри походження, ознаки, пооб’єктні рішення, часові записи, параметри моделей, відомості про середовище та контрольні суми. Незалежна перевірка підтвердила $55$ контрольних сум, без розбіжностей повторно обчислила показники та відтворила $42\,168$ оцінок першого методу з канонічних ознак. Окремий недеструктивний аудит перевірив контрольні суми $2056$ фізичних MP4, повторно сформував ознаки й підтвердив незмінність архівних таблиць; структурні рішення відтворилися, а для спектральної та спільної гілок установлено калібрувальний зсув. Розвідувальні протоколи структурних контрольних прикладів, збереження спектрального сліду й локальних станів технології містять сценарії запуску, пооб’єктні результати, зведені показники та контрольні суми; для останнього додатково збережено $35$ машинозчитуваних записів $S_5$. \textit{Не перевірено.} Повного пакета ретроспективної та первинної внутрішньої конфігурацій немає; первинний спектральний екстрактор і передкодувальні маски не збережено, а нові пакети не містять повних подієво-поведінкових журналів і не відтворюють промисловий п’ятиетапний контур.

\item \textit{Експериментально оцінено у шести відокремлених серіях, окремому ресурсному випробуванні та локальному аудиті технологічних станів.} Ретроспективна перевірка $110$ об’єктів підтвердила арифметичну узгодженість матриці $48$, $7$, $5$, $50$ (у порядку правильно визначених контрольних, хибнопозитивних, хибнонегативних і правильно визначених змінених об’єктів) та похідних показників, а площу під характеристичною кривою не подано через відсутність пооб’єктних оцінок. Внутрішня серія на $400$ об’єктах забезпечила груповий поділ і одноразове оцінювання незмінного порога. Чотири зовнішні набори охопили $38$ природних записів із $12$ пристроїв і $24$ груп походження; до них увійшли перша MP4-серія на $56$ похідних, поетапна серія на $32$, багатоформатна серія на $10\,000$ і відкладена зовнішня серія на $2000$ об’єктах. П’ять форматів --- PNG, JPEG, WebP, MP4 і WebM --- перевірено в багатоформатній та відкладеній зовнішній серіях без повторного навчання першого методу.

У відкладеній зовнішній серії матриця спільної проєкції першого методу становила $424;76;424;1076$, а матриця скороченої моделі узгодження статичних оцінок --- $500;0;508;992$ у порядку правильно визначених контрольних, хибнопозитивних, хибнонегативних і правильно визначених змінених об’єктів. Для моделі узгодження точкова оцінка збалансованої частки правильних рішень змінилася від $0{,}7827$ до $0{,}8307$. Вона усунула $76$ хибнопозитивних рішень, але додала $84$ пропуски, зменшила загальну частку правильних рішень від $0{,}7500$ до $0{,}7460$ і не виявила жодної з $500$ суто спектральних змін. Критерій Мак-Немара дав $p=0{,}5801$, тому встановлено зміну співвідношення між специфічністю та повнотою, а не безумовну перевагу моделі.

Для порівняння конфігурацій обчислено розвідувальний інтегральний рейтинг $K_I$ із чотирьох складників: достовірності виявлення, 95-го процентиля наскрізного часу дослідної послідовності оброблення, найбільшого обсягу фізичної пам’яті процесу та форматної стійкості. Для структурної, спектрально-зорової, спільної й узгоджувальної конфігурацій 95-ті процентилі становили відповідно $2{,}136$, $760{,}662$, $853{,}971$ і $1078{,}904$~мс, пам’ять --- $78{,}77$, $96{,}79$, $101{,}53$ і $110{,}05$~МіБ, а $K_I$ --- $77{,}50$, $50{,}69$, $68{,}95$ і $69{,}46$. Структурна конфігурація мала найбільший інтегральний бал, але не виконала обов’язкової форматної умови через результат $0{,}5000$ на WebM; спільна й узгоджувальна конфігурації цю умову виконали. \textit{Межа висновку.} $K_I$ є вторинним розвідувальним рейтингом, а не ймовірністю правильного рішення; результати підтверджують перевірену статичну основу та скорочене узгодження її оцінок, але не виявлення виконуваного шкідливого коду, повну поведінкову процедуру або промислову експлуатацію.
\end{enumerate}

На формальному рівні сукупність одержаних результатів утворює послідовний підхід до розв’язання поставленої наукової задачі. Ретроспективні агрегати підтримують арифметичні залежності попереднього прогону, внутрішня серія підтверджує діагностичну роздільність детермінованих ознак у межах синтетичного протоколу, а чотири зовнішні набори встановлюють форматні межі статичних конфігурацій на природних записах. Додаткові контрольні аналізи показали, що повне зникнення спектрального сліду не пояснює низького результату $Z_v$, а структурний індикатор не має достатньої специфічності щодо досліджених допустимих розширень контейнера. Відкладена зовнішня серія характеризує перенесення скороченої моделі узгодження статичних оцінок на нові пристрої, ресурсне випробування кількісно описує компроміс між достовірністю, часом, пам’яттю та стійкістю між форматами, а локальний аудит підтверджує порядок $S_1$--$S_5$, атомарну реєстрацію й безпечне завершення за двох класів технічних відмов. Модель визначає, які представлення мультимедійного об’єкта та відомості про їх походження потрібні для аналізу. Метод спектрально-зорового та структурного аналізу формує статичні свідчення (із акцентом на двогілковій інтеграції, а не на самодостатності окремої гілки), а процедура поведінкового узгодження ознак різного походження задає порядок їх причинного зіставлення з доступними подіями оброблення; повну конфігурацію останньої експериментально не підтверджено. Інформаційна технологія перетворює ці результати на простежуваний процес із інваріантами ізоляції, відокремлення $S_A$ від дії політики та реєстрації контексту для наступної перевірки.

Принциповим для запропонованого підходу є явне відображення меж спостереження. Недоступна модальність не трактується як нульове відхилення, порожній подієво-поведінковий журнал не є підтвердженням штатної реакції, а недосягнення порога не доводить абсолютної безпечності носія. У формальній схемі передбачено окреме збереження оцінки ризику, упевненості, локалізацій та дії політики, що дає змогу повторно тлумачити аналітичний результат після зміни правил реагування без підміни початкових свідчень. На рівні формальної конструкції це забезпечує пояснюваність рішення за походженням ознак і водночас обмежує категоричність висновку фактично доступними даними; повну реалізацію цієї властивості не перевірено.

Для детермінованих конфігурацій першого методу виконано групово відокремлений внутрішній поділ, перевірено п’ять форматів, збережено пооб’єктні рішення й часові записи та проведено зовнішнє оцінювання незмінних параметрів. Скорочену модель узгодження статичних оцінок окремо перевірено на $2000$ об’єктах із чотирьох нових пристроїв, однак результат визначався лише статичними оцінками й не замінює повного поведінкового узгодження. Локальна перевірка переходів і відмов також не замінює випробування розподіленої системи. Подальше підтвердження повної конфігурації потребує нових наперед відокремлених природних джерел, причинно пов’язаних послідовностей подієво-поведінкових відомостей, випробування інших форматів і кодеків, незалежного повторення ресурсного оцінювання, реконструкції, зовнішніх дій і тривалого промислового навантаження. Ці умови не змінюють доказового статусу ретроспективних агрегатів і не дозволяють переносити показники на невипробувані вебсистеми.

\clearpage
